\section{特性方程式と漸化式}\label{節：特性方程式と漸化式}
\subsection{不動点からのずれを追う}

まず
\[
 a_1=1,\qquad a_{n+1}=2a_n+1
\]
を考える. 項は$1,3,7,15,\ldots$と進む.
毎回1を足すため,\,元の数列は等比数列ではない.
しかし,\,各項に1を加えると
\[
 a_{n+1}+1=2(a_n+1)
\]
となる. $b_n=a_n+1$は初項2,\,公比2の等比数列だから
\[
 b_n=2^n,\qquad a_n=2^n-1.
\]

なぜ1を加えるのだろうか.
更新$f(x)=2x+1$で変わらない値を探すと,
$\alpha=2\alpha+1$から$\alpha=-1$となる.
このような$f(\alpha)=\alpha$を満たす値を\textbf{不動点}という.
先ほど追った$a_n+1$は,\,不動点からのずれ$a_n-\alpha$であった.

一般に
\[
 a_{n+1}=pa_n+q,\qquad p\ne1
\]
では,\,$\alpha=\frac{q}{1-p}$が不動点になる.
二つの式を引けば
\[
 a_{n+1}-\alpha=p(a_n-\alpha).
\]
従って
\[
 a_n=\alpha+p^{n-1}(a_1-\alpha).
\]
ただし$p=0$では,\,$n=1$は初項そのものを読み,
$n\ge2$では$a_n=q$とする.
零の零乗を計算する必要はない.

\NewFigFixed
原点を$\alpha$へ移せば,\,毎回の更新は$p$倍だけになる.
$p<0$なら向きが反転する.
$|p|<1$ならずれは零に近づき,\,数列は$\alpha$へ収束する.
$|p|>1$では,\,不動点から出発する場合を除いてずれの絶対値が増大する.
$p=-1$なら,\,不動点を挟んだ二つの値を交互に取る.
式変形は,\,数直線上の動きも説明していたのである.

\subsection{定数項も一つの成分にする}

この更新を行列で書くため,\,変化しない数1をもう一つの成分に加える.
\[
 \boldsymbol x_n=\begin{pmatrix}a_n\\1\end{pmatrix},
 \qquad
 A=\begin{pmatrix}p&q\\0&1\end{pmatrix}
\]
とすれば
\[
 \boldsymbol x_{n+1}=A\boldsymbol x_n.
\]
下の成分は常に1で,\,上の成分だけが元の規則に従う.
記録する成分を一つ増やすと,\,定数の足し算も行列の積に含められる.

$p\ne1$では
\[
 \begin{pmatrix}a_n\\1\end{pmatrix}
 =\begin{pmatrix}\alpha\\1\end{pmatrix}
 +(a_n-\alpha)\begin{pmatrix}1\\0\end{pmatrix}.
\]
前者は更新しても変わらず,\,後者の方向は$p$倍になる.
つまり,\,不動点とずれを分ける解法は,\,\textbf{変化の倍率が異なる二方向への分解}だった.

$p=1$の場合も行列の表示は使える.
\[
 \begin{pmatrix}1&q\\0&1\end{pmatrix}^{m}
 =\begin{pmatrix}1&mq\\0&1\end{pmatrix}
\]
なので,\,$a_n=a_1+(n-1)q$となる.
$q\ne0$では不動点がなく,\,ずれの倍率だけには直せない.
足し算が積み重なって$n$が現れることを覚えておこう.

\subsection{2階の式を1階の更新へ}

次は
\[
 a_{n+2}=5a_{n+1}-6a_n
\]
を考える.
$a_n$一つだけでは次の状態を決められないが,\,隣り合う二項を持てばよい.
\[
 \boldsymbol v_n=\begin{pmatrix}a_{n+1}\\a_n\end{pmatrix},
 \qquad B=\begin{pmatrix}5&-6\\1&0\end{pmatrix}
\]
とすると
\[
 B\boldsymbol v_n
 =\begin{pmatrix}5a_{n+1}-6a_n\\a_{n+1}\end{pmatrix}
 =\boldsymbol v_{n+1}.
\]
第二行は,\,上の成分を一つ古い記録として下へ移す.
過去の二項が必要な式も,\,\textbf{二成分の状態については1階の更新}である.

そこで,\,掛けても向きが変わらず,\,定数倍になるベクトルを探す.
この例では
\[
 B\begin{pmatrix}2\\1\end{pmatrix}
 =2\begin{pmatrix}2\\1\end{pmatrix},
 \qquad
 B\begin{pmatrix}3\\1\end{pmatrix}
 =3\begin{pmatrix}3\\1\end{pmatrix}.
\]
二方向に分ければ,\,一方は2倍,\,他方は3倍されるだけになる.

一般に零でない$\boldsymbol u$が
\[
 B\boldsymbol u=\lambda\boldsymbol u
\]
を満たすとき,\,$\boldsymbol u$を\textbf{固有ベクトル},
$\lambda$を対応する\textbf{固有値}という.
固有値は,\,その方向の一回の更新の倍率である.
零ベクトルはどんな$\lambda$でも式を満たすため,\,定義から除く.

\subsection{特性方程式は倍率を決める}\label{節：倍率を決める方程式}

一般の2階の式
\[
 a_{n+2}=s a_{n+1}+t a_n
\]
では,\,$B=\begin{pmatrix}s&t\\1&0\end{pmatrix}$である.
固有ベクトルを$\begin{pmatrix}x\\y\end{pmatrix}$とすると
\[
 sx+ty=\lambda x,\qquad x=\lambda y.
\]
$y=0$なら$x=0$にもなるので,\,固有ベクトルでは$y\ne0$である.
$y$で割って$\begin{pmatrix}\lambda\\1\end{pmatrix}$にそろえると,
\[
 s\lambda+t=\lambda^2
\]
が必要になる. 従って,\,倍率は
\[
 \lambda^2-s\lambda-t=0
\]
という\textbf{特性方程式}から決まる.
その根を\textbf{特性根}とも呼ぶ.
これが,\,漸化式の解法に二次方程式が現れる理由である.

二行二列の一般の行列
$A=\begin{pmatrix}p&q\\r&s\end{pmatrix}$
でも,\,$(\lambda I-A)\boldsymbol u=\boldsymbol0$に零でない解がある条件は
\[
 \begin{aligned}
 \det(\lambda I-A)
 &=(\lambda-p)(\lambda-s)-qr\\
 &=\lambda^2-(p+s)\lambda+ps-qr=0.
 \end{aligned}
\]
前節の行列式によって,\,固有値を求められる.

\begin{bookpoint}
$\alpha=f(\alpha)$は\textbf{変わらない値}を探す.
$A\boldsymbol u=\lambda\boldsymbol u$は\textbf{一定倍率で変わる方向}を探す.
不動点の値$\alpha$と倍率$\lambda$は,\,同じ役割の数ではない.
\end{bookpoint}

\subsection{二つの等比数列に分ける}

二つの異なる根を$\lambda_1,\lambda_2$とする.
対応する$\begin{pmatrix}\lambda_1\\1\end{pmatrix}$と
$\begin{pmatrix}\lambda_2\\1\end{pmatrix}$は比例しないので,\,初期状態を一意に
\[
 \begin{pmatrix}a_2\\a_1\end{pmatrix}
 =c\begin{pmatrix}\lambda_1\\1\end{pmatrix}
 +d\begin{pmatrix}\lambda_2\\1\end{pmatrix}
\]
と分けられる.
各方向を$n-1$回更新して下の成分を読むと,
\[
 a_n=c\lambda_1^{n-1}+d\lambda_2^{n-1}.
\]
根が非零なら通常の冪として読み,\,根が零ならその成分は初期状態にだけ寄与し,\,一回の更新で消えると読む.
初期値を合わせた後は更新の一意性によって,\,全項がこの式に一致する.

\noindent\textbf{【例題】}\par\nobreak
$a_1=1,\,a_2=4$,\,$a_{n+2}=5a_{n+1}-6a_n$の一般項を求めよ.

\noindent\textbf{【方針】}\par\nobreak
2倍と3倍の方向に初期状態を分けます.

\noindent\textbf{【解答】}\par\nobreak
\[
 a_n=-2^{n-1}+2\cdot3^{n-1}.
\]
\noindent\textbf{【解法】}\par\nobreak
特性根は2,\,3である. 初期状態は
\[
 \begin{pmatrix}4\\1\end{pmatrix}
 =-\begin{pmatrix}2\\1\end{pmatrix}
 +2\begin{pmatrix}3\\1\end{pmatrix}.
\]
従って
\[
 \boldsymbol v_n
 =-2^{n-1}\begin{pmatrix}2\\1\end{pmatrix}
 +2\cdot3^{n-1}\begin{pmatrix}3\\1\end{pmatrix}.
\]
下の成分から一般項を得る.
$n=1,2$で初期値に合い,\,二つの冪はそれぞれ更新式を満たす.
線形性によって,\,その和も更新式を満たす.

\subsection{補助数列を取り出す}

同じ式は
\[
 a_{n+2}-2a_{n+1}
 =3(a_{n+1}-2a_n)
\]
と変形しても解ける.
この$a_{n+1}-2a_n$は,\,行列のどこにあるのだろうか.

数を横に並べた$L=(u,v)$を\textbf{行ベクトル}と呼び,
$L\begin{pmatrix}x\\y\end{pmatrix}=ux+vy$と定める.
先ほどの$B$について
\[
 (1,-2)B=(3,-6)=3(1,-2)
\]
だから
\[
 \begin{aligned}
 b_n&=(1,-2)\boldsymbol v_n=a_{n+1}-2a_n,\\
 b_{n+1}&=(1,-2)B\boldsymbol v_n=3b_n.
 \end{aligned}
\]
このように$LA=\lambda L$を満たす零でない行ベクトルを\textbf{左固有ベクトル}という.

固有ベクトルは,\,一定倍率になる\textbf{状態の方向}を見つける.
左固有ベクトルは,\,一定倍率になる\textbf{取り出し方}を見つける.
二つの解法は,\,同じ更新を別の側から見ている.

\subsection{重解で現れる足し算}\label{節：重解と足し算}

特性方程式が重解$r\ne0$を持つと
\[
 a_{n+2}=2r a_{n+1}-r^2a_n.
\]
$a_n=r^{n-1}b_n$と置いて倍率を取り除けば
\[
 b_{n+2}-2b_{n+1}+b_n=0.
\]
これは$b_{n+2}-b_{n+1}=b_{n+1}-b_n$である.
差が一定なので
\[
 b_n=C+D(n-1),
\]
従って
\[
 a_n=r^{n-1}\{C+D(n-1)\}.
\]
等比的な変化を取り除いた後に,\,等差的な変化が残った.
重解だから同じ冪を二回書くのでは,\,二つの初期値を自由に合わせられない.

行列でも,\,この足し算を直接見られる.
例えば$r=2$では
\[
 B=\begin{pmatrix}4&-4\\1&0\end{pmatrix},
 \quad
 \boldsymbol u=\begin{pmatrix}2\\1\end{pmatrix},
 \quad
 \boldsymbol w=\begin{pmatrix}1\\0\end{pmatrix}
\]
とすると
\[
 B\boldsymbol u=2\boldsymbol u,\qquad
 B\boldsymbol w=\boldsymbol u+2\boldsymbol w.
\]
$\boldsymbol w$は固有ベクトルではないが,\,更新すると$\boldsymbol u$の方向へも寄与する.
このように$(B-rI)\boldsymbol w=\boldsymbol u$,
$(B-rI)\boldsymbol u=\boldsymbol0$となる$\boldsymbol w$は,\,\textbf{一般化固有ベクトル}の例である.

$\boldsymbol v_n=c_n\boldsymbol u+d_n\boldsymbol w$と書くと
\[
 d_{n+1}=2d_n,\qquad c_{n+1}=2c_n+d_n.
\]
各係数から$2^{n-1}$を取り除けば,\,第二の式は一定量を足す式になる.
毎回の寄与が積み重なるため,\,$n$が現れるのである.

\noindent\textbf{【例題】}\par\nobreak
$a_1=1,\,a_2=4$,\,$a_{n+2}=4a_{n+1}-4a_n$の一般項を求めよ.

\noindent\textbf{【方針】}\par\nobreak
公比2の成長を取り除き,\,差を調べます.

\noindent\textbf{【解答】}\par\nobreak
\[
 a_n=n2^{n-1}.
\]
\noindent\textbf{【解法】}\par\nobreak
$b_n=\frac{a_n}{2^{n-1}}$とすると
$b_1=1,\,b_2=2$で,\,二階差が零になる.
従って$b_n=n$である.
行列では$\boldsymbol v_1=\boldsymbol u+2\boldsymbol w$に対応する.

なお,\,一般の行列は固有値が重なっても,\,固有ベクトルを十分に選べる場合がある.
例えば$2I$はどの方向も2倍にする.
この場合は足し算の蓄積がなく,\,すべての成分が単に2倍される.
\textbf{重解そのものと,\,固有方向が足りないことは区別する}.
$r=0$の2階の式は$a_{n+2}=0$なので,\,最初の二項だけを別に記録すればよい.
